\FloatBarrier

\section{Experimental details}
\label{sec:details}

\begin{table}[t]
\centering
\footnotesize
\setlength{\tabcolsep}{3.4pt}
\begin{tabular}{lcccccc}
\toprule
& \multicolumn{3}{c}{Latency (s)} & \multicolumn{2}{c}{Speedup} & \\
\cmidrule(lr){2-4}\cmidrule(lr){5-6}
Date & AR & \vispec & \hivis & \vispec & \hivis & TTFT \% \\
\midrule
\csname @@input\endcsname generated/jobs_single.tex
\bottomrule
\end{tabular}
\caption{\textbf{Every single-image job.} TTFT \% is the share of the time to the first token in the latency of \vispec.}
\label{tab:jobs-single}
\end{table}

\noindent\textbf{Prompts.}
Every Qwen prompt begins with the same system turn, ``A chat between a curious human and an artificial intelligence assistant. The assistant gives helpful, detailed, and polite answers to the human's questions.''
\Cref{tab:prompts} lists the instruction of each benchmark and style.
On single images the image follows the instruction, and in a manga window the pages come first, in reading order, followed by the instruction.

\begin{table}[h]
\centering
\footnotesize
\begin{tabular}{p{0.24\linewidth}p{0.66\linewidth}}
\toprule
Benchmark and style & Instruction \\
\midrule
TextVQA, long & Perform an OCR task on the provided image. Please extract the text accurately and provide a detailed explanation of the process. Ensure the response is comprehensive and well-structured. \\
TextVQA, short & Answer the question using a single word or phrase. \\
COCO, long & Please provide a detailed description of the given image. \\
COCO, short & Provide a one-sentence caption for the provided image. \\
Manga window & Summarize the story told in these manga pages, in order. \\
\bottomrule
\end{tabular}
\caption{\textbf{Instructions of the Qwen prompts.} TextVQA prompts also carry the question.}
\label{tab:prompts}
\end{table}

\noindent\textbf{Image resolution.}
Single images are limited to $1280 \times 28 \times 28$ pixels.
Manga pages are limited to $1200 \times 28 \times 28$ pixels, so that the longest window, the 1,024 new tokens and the draft tree fit within the 8,192 positions of the \vispec draft model's rotary table for \qwenS.
Each page passes through the vision encoder in its own call.

\noindent\textbf{Cache and memory.}
Every Qwen run sizes the KV cache for each sample, to hold the kept prompt, the new tokens and the draft tree.
Live KV is the arithmetic size of that cache at the end of a sample.

\noindent\textbf{Software.}
All Qwen runs use PyTorch 2.7.0 with CUDA 12.8, Transformers 4.51.3 and the released checkpoints at fixed revisions.

\noindent\textbf{Runs.}
\Cref{tab:jobs-single,tab:jobs-pages} list every Qwen job whose timings enter the report, one row for each benchmark, style and job, with the date it ran.
Each job ran on one pinned GPU, and the outputs of each method are token-identical across the jobs of one row group.

\FloatBarrier
\section{Results of every job}
\label{sec:jobs}



\begin{table}[h]
\centering
\footnotesize
\setlength{\tabcolsep}{2.6pt}
\begin{tabular}{llcccccc}
\toprule
& & \multicolumn{3}{c}{Latency (s)} & & \multicolumn{2}{c}{Speedup} \\
\cmidrule(lr){3-5}\cmidrule(lr){7-8}
Target & Date & AR & \vispec & \hivis & TTFT (s) & \vispec & \hivis \\
\midrule
\csname @@input\endcsname generated/jobs_pages.tex
\bottomrule
\end{tabular}
\caption{\textbf{Every manga-window job.} TTFT is the mean time to the first token of AR.}
\label{tab:jobs-pages}
\end{table}

The \qwenS manga jobs show the largest spread between jobs with identical outputs.
The AR latency of the slowest job exceeds that of the fastest by \tbPagesArSpread\%, while the speedup stays within \tbPagesSpeedMin to \tbPagesSpeedMax.
On \qwenL the AR latency varies by \sbPagesArSpread\% across the three jobs, which ran the three methods in rotated order, and the lead of \hivis over \vispec stays at \sbPagesHovMin.

\FloatBarrier
\section{Exactness checks}
\label{sec:exact-details}

\noindent\textbf{Full-precision replays.}
On \qwenS the replays cover the first five long TextVQA questions under three eviction settings, \fpqTbSingle samples, and the first three manga windows with the full cache, with KV eviction at a quarter and with both forms at a quarter, \fpqTbPages samples.
On \qwenL they cover \fpqSbSingle single-image samples of \vispec and \hivis in both image positions and \fpqSbPages manga-window samples of the two drafters.
Every replay is identical to AR.

\noindent\textbf{Divergence diagnosis.}
The diagnosis covers \diagCasesTb first differences of \vispec on \qwenS, from long and short single-image answers and from the manga windows, \diagCasesSb of \vispec and \hivis on single images of \qwenL and \diagCasesSbPages on the manga windows of \qwenL.
All are near-ties except the two differences on the manga windows of \qwenL described in \cref{sec:exactness}.

\noindent\textbf{\hivis on several images.}
With one image the multi-image mask hands the prompt to the released mask, and five long TextVQA answers of \hivis on \qwenL are token-identical to those of the single-image runs.
On the manga windows the full-precision replay of \hivis matches AR on all three windows, which also covers the position bookkeeping of the removed pages.
